\chapter{案例设计与分析方法}
\section{职业阶段与材料组织}
依据媒介重心与职业定位的变化，设定时尚模特、影视转型、多领域拓展、职业稳定与新媒体重塑五个阶段。为使每条记录只归属一个阶段，本文采用互不重叠的年份区间，见表\ref{tab:career-stages}。这些区间是模拟编码规则，不是经核验的真实履历。

\begin{table}[htbp]
  \centering
  \caption{虚构案例的职业阶段与观察维度}
  \label{tab:career-stages}
  \small
  \begin{tabular}{llll}
    \toprule
    阶段 & 时间区间 & 设定的主要媒介 & 观察维度 \\
    \midrule
    时尚模特期 & 2002--2007 & 杂志与平面摄影 & 造型与视觉风格 \\
    影视转型期 & 2008--2009 & 影视与宣传报道 & 角色与职业定位 \\
    多领域拓展期 & 2010--2016 & 影视与综艺 & 媒介组合 \\
    职业稳定期 & 2017--2019 & 综合媒体 & 形象一致性 \\
    新媒体重塑期 & 2020--2025 & 视频与社交平台 & 互动与内容更新 \\
    \bottomrule
  \end{tabular}
\end{table}

每条记录包含年份、媒介类型、阶段编号与内容标签。图片记录保留画面形式与使用场景，文本记录保留主题与情感类别。本文选取两张人物图片作为视觉观察材料，分别记录服饰、姿态、构图和背景；由于缺少经核验的拍摄时间与媒介来源，不将这两张图片用于判定职业阶段。

\section{研究流程}
分析按材料整理、阶段编码、指标计算与结果解释四步进行，见图\ref{fig:workflow}。先统一记录格式，再划分阶段与时间窗口，随后计算情感得分并复核编码一致性，最后结合情境设定解释差异。
\begin{figure}[htbp]
  \centering
  \begin{tikzpicture}[>=stealth]
    \node[draw, rounded corners, minimum width=3.8cm, minimum height=0.8cm] (a) at (0,0) {整理图片与模拟文本};
    \node[draw, rounded corners, minimum width=3.8cm, minimum height=0.8cm] (b) at (0,-1.3) {划分阶段与编码材料};
    \node[draw, rounded corners, minimum width=3.8cm, minimum height=0.8cm] (c) at (0,-2.6) {计算情感与一致性指标};
    \node[draw, rounded corners, minimum width=3.8cm, minimum height=0.8cm] (d) at (0,-3.9) {比较结果并说明边界};
    \draw[->,thick] (a) -- (b);
    \draw[->,thick] (b) -- (c);
    \draw[->,thick] (c) -- (d);
  \end{tikzpicture}
  \caption{虚构职业发展案例的分析流程}
  \label{fig:workflow}
\end{figure}

\section{情感得分与编码一致性}
将评论编码为正向、中性与负向三类，分别赋值为 $+1$、$0$ 与 $-1$。时间窗口 $t$ 的加权情感得分定义为
\begin{equation}
  S(t)=\frac{\sum_{i=1}^{N_t}w_i s_i(t)}{\sum_{i=1}^{N_t}w_i},
  \qquad w_i\geq 0,\quad \sum_{i=1}^{N_t}w_i>0.
  \label{eq:sentiment}
\end{equation}
其中，$N_t$ 为窗口内评论数，$w_i$ 为评论权重，$s_i(t)$ 为情感编码。式\eqref{eq:sentiment}将得分限制在 $[-1,1]$，避免权重规模改变得分范围。为便于复算，本例全部取 $w_i=1$，于是
\begin{equation}
  S(t)=\frac{n_{+}(t)-n_{-}(t)}{N_t},
  \qquad N_t=n_{+}(t)+n_{0}(t)+n_{-}(t).
  \label{eq:sentiment-count}
\end{equation}

两位编码者对同一批模拟文本独立判断，用观察一致率 $p_o$ 与偶然一致率 $p_e$ 计算
\begin{equation}
  \kappa=\frac{p_o-p_e}{1-p_e}.
  \label{eq:kappa}
\end{equation}
本例设定 $p_o=0.90$、$p_e=1/3$，则 $\kappa=0.85$。这只是公式计算示例；实际研究应由双方完整编码结果计算 $p_o$ 与 $p_e$。
